% ----------------------------------------------------------------------------
%                               Conclusão
% ----------------------------------------------------------------------------
\chapter{Conclusão}
\label{chap:conclusão}

O presente trabalho teve como objetivo central investigar e propor uma solução tecnológica voltada ao suporte pedagógico e ao manejo comportamental de alunos com Transtorno do Déficit de Atenção com Hiperatividade (TDAH) e Transtorno Opositivo Desafiador (TOD). Por meio do levantamento de requisitos e da subsequente prototipagem de alta fidelidade, foi possível materializar uma aplicação web responsiva que visa atender às demandas latentes da comunidade escolar.

A principal contribuição deste projeto reside na proposição de uma transição do modelo de registro passivo para uma abordagem de manejo proativo. A interface desenvolvida, estruturada sob os princípios de usabilidade e baixa carga cognitiva, apresenta-se como uma alternativa viável para o dinâmico contexto escolar. A substituição de formulários extensos por seleções ágeis (chips) e a entrega imediata de estratégias de intervenção baseadas em evidências têm o potencial de reduzir significativamente a sobrecarga do corpo docente em momentos de crise.

Paralelamente, a estruturação de uma linha do tempo (timeline) dedicada aos familiares democratiza o acesso à informação, promovendo uma comunicação transparente e livre de ruídos entre a escola e o ambiente doméstico. Essa arquitetura reafirma a importância da tecnologia não apenas como ferramenta de arquivamento, mas como ponte para a construção de uma rede de apoio colaborativa em prol do desenvolvimento do estudante neurodivergente.

Devido às restrições temporais inerentes ao escopo deste trabalho, a etapa de validação prática da ferramenta em ambiente escolar e a aplicação de testes de usabilidade com o público-alvo não puderam ser realizadas. Como trabalhos futuros, sugere-se a implementação do código-fonte da aplicação web com base nos requisitos levantados, seguida de um estudo de caso prático no Instituto Federal de Brasília (IFB), visando avaliar o impacto real e a aplicabilidade dos cards de manejo proativo na rotina dos educadores e das famílias.